%!TEX root = ../main.tex
% =============================================================================
%  十、模型评价与推广            建议 2 页
%  评审在问：你诚实吗？能推广吗？
%  本节必须出现：① 局限\textbf{四要素}（指名假设 + 影响方向 + 数据需求 + 校准动作）
%                ② 推广的\textbf{临界条件}  ③ 行业口径声明（电能量收益的保守下界）
%                ④ 创新五段式（缺陷—改动—证据—代价—边界）
%  禁止空话："数据有限""模型有待改进"。
% =============================================================================

\section{模型评价与推广}

\subsection{主要创新点}

\begin{guideblock}
每个创新都用\textbf{五段式}写，缺一段就不算创新：
\textbf{缺陷（用数据/残差/约束分析找证据）\(\rightarrow\) 改动（必须能写成公式）
\(\rightarrow\) 证据（来自消融矩阵，同一场景集）\(\rightarrow\) 代价 \(\rightarrow\) 边界}。
禁止写成方法罗列（"本文采用了 X 方法"）；禁止出现"显著提升"等无对照表述。
"能说出自己的方法何时失效，是最高级的自信。"
\end{guideblock}

\begin{innovation}{1}{基于对偶水价值的阈值型充放电策略}
\ip{缺陷}直接对 LP 求解只能得到一个数值解，无法解释"为什么在这些时段充电、
在那些时段放电"，也无法回答"容量摆幅与充电功率是否仍是瓶颈"（放电功率不紧）；当参数变化时需要重新求解才能给出判断。
\ip{改动}本文由 LP 对偶得到储电的边际价值（水价值）\(\mu_t\)，
把最优性条件整理为阈值律（命题~\ref{prop:threshold}）：
当 \(p_t\le\eta_c\mu_{t}\) 时充电，当 \(p_t\ge\mu_{t}/\eta_d\) 时放电，
中间构成死区（式~\eqref{eq:q1-deadband}）。与基线的差别在于：
阈值律给出了\textbf{与求解器无关}的策略形式，且把"是否值得套利"归结为
价比阈值 \(1/(\eta_c\eta_d)\approx1.235\)。
\ip{证据}在纯套利子问题/储电量界不活跃时，阈值律一致性检验的违反次数为 \ph，
两条恒等式的最大偏差为 \ph；由水价值曲线（图~\ref{fig:watervalue}）读出
储电量的边际价值归零点为 \ph~kWh，据此判断容量摆幅是否为瓶颈，并另报充电功率的边际价值（见 \evid{C-Q1-02}）。
\ip{代价}需要求解并保存对偶变量，计算量增加 \ph；对偶解释需要额外的符号与推导说明。
当储电量上下界活跃时，阈值形式失效，须退回互补松弛形式
式~\eqref{eq:q1-kkt-c}--\eqref{eq:q1-kkt-d}。
\ip{边界}当 \(\overline P\) 收紧至 \ph~kW 以下、充电功率成为主导瓶颈时，
阈值律需扩展为"功率约束下的分段阈值"，此时应退回完整的互补松弛形式。
\end{innovation}

\begin{innovation}{2}{以预报频次对照实验量化信息价值（VoI）}
\ip{缺陷}题面直接追问"是否需要引入其他时刻的预报"，而基线做法只能定性回答
"建议引入更多预报"，没有把信息折算成钱。
\ip{改动}本文设计预报时刻数 \(n\in\{1,2,4,8,24\}\) 的对照实验，
在同一固定场景集上重解模型，定义信息价值
\(\mathrm{VoI}(n)=\bigl(C(1)-C(n)\bigr)/C(1)\)（式~\eqref{eq:q3-voi}），
并同时给出\textbf{成本侧}估计。
\ip{证据}总费用在 \(n=\ph\) 之后进入 \(\pm0.5\%\) 平稳带；
由 \(n=4\) 加密到 \(n=\ph\) 仅再节省 \ph~元（占总费用 \ph），
低于估算的加密成本 \ph~元（图~\ref{fig:voi}）（见 \evid{C-Q3-01}）。
\ip{代价}需要额外求解 \ph 组模型，计算时间增加 \ph；
加密成本估计依赖外部报价，属于估计而非本队实测。
\ip{边界}该结论依赖于本题的预报误差水平与惩罚倍数；
当紧急购电倍数 \(\kappa\) 或预报误差放大到 \ph 倍以上时，平稳带右移，
需重新计算后方可套用。
\end{innovation}

\begin{guideblock}
上两条是模板给出的示例写法，请按本队实际工作替换或增补。
验收标准：每个创新都有\textbf{消融证据}；"边界"段非空；
不出现"显著提升"等无对照表述；不出现无法核实的行业数字。
\end{guideblock}

\subsection{模型优点}

\begin{itemize}
  \item \textbf{结构统一}：四问共用同一套物理模型骨架（式~\eqref{eq:q1-balance}--
        \eqref{eq:q1-nonneg}），差别只在信息结构与风险处理，
        因此每一层的新增贡献都可被单独消融，避免了"四篇互不相干报告"的组织方式；
  \item \textbf{可解释}：对偶水价值与阈值律给出与求解器无关的策略形式，
        使"为什么这样充放"可以被解释和被质疑；
  \item \textbf{可复算}：全部数值均由一键脚本从原始附件重算，
        并通过原始可行性、对偶最优性与费用独立复算三重检验。
\end{itemize}

\subsection{模型局限}

\begin{guideblock}
局限必须写\textbf{四要素}，缺一不可：
① 哪条假设可能偏离现实（指名编号）；
② 偏离会让结论往哪个方向变（给符号：高估还是低估）；
③ 什么数据能验证（给具体数据需求）；
④ 部署前如何校准（给具体动作）。
\end{guideblock}

\begin{table}[htbp]
  \centering
  \caption{模型局限的四要素表述}
  \label{tab:limitation}
  \footnotesize
  \begin{tabularx}{\linewidth}{C{1.6cm} L{3.6cm} C{1.8cm} L{3.2cm} Y}
    \toprule
    假设编号 & 可能偏离现实之处 & 偏离方向 & 验证所需数据 & 部署前校准动作 \\
    \midrule
    A2 & 充放电效率取单侧 0.9，实际与功率、温度、老化相关
       & 高估收益 & 电池厂商的效率—功率—温度曲线、运行日志
       & 按 \(\eta\) 灵敏度结果重新标定；若 \(\eta_c\eta_d<0.7\) 则套利收益需整体下调 \\
    A4 & 忽略自放电与循环衰减
       & 高估长期收益 & 循环寿命试验数据、衰减—循环次数曲线
       & 在目标函数中加入循环成本 \(c_{cyc}\)，按衰减灵敏度重新优化 \\
    A1 & 不允许向外网售电
       & 高估费用（结果保守） & 上网电价政策与结算规则
       & 引入上网电价 \(\theta\) 后重解，比较策略变化 \\
    R3 & 采用日前预报而非完美信息
       & 高估紧急购电费用 & 实测预报误差序列（更长历史）
       & 按预报误差 \(\times0.5/\times1/\times2\) 灵敏度结果插入安全边际 \\
    \bottomrule
  \end{tabularx}
\end{table}

本文目标函数只含电能量购电费用，未计入容量补偿、辅助服务收益与绿电环境价值，
也未计入储能循环衰减与安全成本；因此\textbf{本文所得节约幅度应视为储能电能量价值的保守下界}。

\begin{guideblock}
上面这一句是\textbf{行业口径必须写进去的一句}，请原样保留。
它的作用是主动声明结论的边界，避免被质疑"把储能收益吹大了"。
\end{guideblock}

\subsection{模型推广与临界条件}

\begin{guideblock}
推广必须给\textbf{临界条件}，而不是"本文模型可以推广到其他场景"这种空话。
同时\textbf{不得偷偷改动题设}：负载仍刚性、不外送、储能收益仍只有价差。
允许在推广中讨论的内容：容量电价与多元收益、系统运行费等附加电费、
虚拟电厂与需求响应、售电与绿电环境价值——但必须放在本节，不得进入主模型。
\end{guideblock}

本文模型可推广到以下情形，各自附带成立条件：

\begin{itemize}
  \item \textbf{峰谷价差更小的地区}：当谷峰价比低于盈亏平衡阈值
        \(1/(\eta_c\eta_d)\approx1.235\) 时，纯套利不再有利，
        模型退化为"以规避紧急购电为主"的保险型策略，此时 \(C^{base}\) 与最优费用的差距迅速缩小；
  \item \textbf{允许余电上网}：需在式~\eqref{eq:q1-balance} 中引入上网电量变量
        并在目标中加入售电收益 \(\theta_t\)，弃光变量 \(s_t\) 的地位由"必要"变为"可选"；
        当 \(\theta_t>0\) 时，最优解中 \(s_t\) 通常为 0；
  \item \textbf{引入需求响应}：负载变为可调后，式~\eqref{eq:q1-balance} 的 \(L_t\) 成为决策量，
        模型由 LP 变为含柔性负荷约束的 LP，本文的阈值律结论需重新推导
        （水价值与负荷转移的边际价值共同决定阈值）；
  \item \textbf{计入附加电费}：用户到户电价 = 电能量费 + 输配电费 + 系统运行费 + 政府性基金及附加。
        若附加电费按电量比例收取，则等于对所有电价整体加价，最优策略不变；
        若按容量（需量）收取，则会显著改变策略，需在目标函数中增加需量费用项后重解。
\end{itemize}

\begin{guideblock}
上一条"若附加电费按电量比例收取则最优策略不变"是一个\textbf{可证明的结论}
（线性目标下所有电价同乘一个因子不改变最优解），
比空泛地"考虑政策影响"有价值得多。若本队有余力，可以在此补一句证明思路。
\end{guideblock}
